\section{Key Concepts}

\begin{keybox}{1\enspace Scaled Dot-Product Attention}
Given queries \(Q\), keys \(K\), and values \(V\) (each a matrix of row-vectors),
attention is:
\[
  \operatorname{Attention}(Q,K,V)
  = \operatorname{softmax}\!\left(\frac{QK^{\top}}{\sqrt{d_k}}\right)V
\]
The \(\sqrt{d_k}\) scaling prevents vanishingly small softmax gradients when the
embedding dimension is large.
\end{keybox}

\vspace{5pt}

\begin{keybox}{2\enspace Multi-Head Attention (MHA)}
Rather than a single attention function, MHA runs \(h\) heads in parallel on
learned linear projections, then concatenates and re-projects:
\[
  \operatorname{MHA}(Q,K,V)
  = \operatorname{Concat}(\mathrm{head}_1,\ldots,\mathrm{head}_h)\,W^{O}
\]
Different heads capture complementary relational patterns: syntactic agreement,
coreference, positional proximity, etc.
\end{keybox}

\vspace{5pt}

\begin{keybox}{3\enspace Positional Encoding}
Attention is permutation-invariant, so position is injected via sinusoidal signals
added to the input embeddings:
\[
  \mathrm{PE}_{(p,2i)}   = \sin\!\bigl(p/10000^{2i/d}\bigr),\qquad
  \mathrm{PE}_{(p,2i+1)} = \cos\!\bigl(p/10000^{2i/d}\bigr)
\]
Modern variants (RoPE, ALiBi) modify attention scores directly for better
length generalisation.
\end{keybox}

% ─── Worked Example / Diagram ─────────────────────────────
